\setcounter{equation}0
\renewcommand{\t}[1]{\textrm{#1}}
\newcommand{\ket}[1]{|#1\rangle}
\newcommand{\bra}[1]{\langle#1|}
\newcommand{\braket}[2]{\langle #1 | #2 \rangle}

\vskip2cm
\advance\baselineskip by.2pt
\spis{Co nieco o~energii}
{Ludwik Lehman}

\wtyt{Co nieco o~energii\hfill
\waut{Ludwik LEHMAN}}
Energia to kluczowe pojęcie nie tylko w~fizyce, lecz także w~innych naukach przyrodniczych. Jednak bardzo trudno znaleźć w~miarę precyzyjną definicję energii. Z~jednej strony to nie dziwi. Podstawowe pojęcia są do zdefiniowania najtrudniejsze. \emph{Co to jest czas? Kiedy nikt mnie nie pyta, wiem. Ale kiedy chcę wytłumaczyć komuś, kto o~to pyta, już nie wiem.}
To stwierdzenie św. Augustyna dobrze pokazuje, na czym polega dylemat operowania podstawowymi pojęciami. Jednak patrząc z~drugiej strony, potrzebujemy odpowiednich nazw oraz ich choćby roboczych, intuicyjnych definicji. \emph{W~fizyce nazwy pomagają wytworzyć stan umysłu, w~jakim postrzegamy pojęcie fizyczne. Dobra nazwa wywołuje w~umyśle obraz, który uwypukla najważniejsze własności pojęcia i~tym samym, w~sposób podświadomy, intuicyjny pomaga zapoczątkować owocne badania. Zła nazwa może wytworzyć blokadę umysłową, która przeszkadza w~prowadzeniu badań.} Tak trafnie napisał Kip S. Thorne w~książce ,,Czarne dziury i~krzywizny czasu''. Gdy zamienimy w~tym cytacie słowo ,,nazwa'' na ,,definicja pojęcia'' – otrzymamy równie trafne spostrzeżenie.

\marg[20]{
\includegraphics{art/202501_kol17}}%
Podręczniki akademickie raczej unikają mierzenia się z~,,intuicyjną'' definicją energii, zaczynając od operacyjnych definicji energii kinetycznej, potencjalnej etc. Na przykład tak robią Kittel, Knight i~Ruderman w~,,Mechanice'' (część tzw. kursu berkelejowskiego). Słynne ,,Feynmana wykłady z~fizyki'' w~podrozdziale \emph{Co to jest energia?} oferują historyjkę o~Piekielnym Piotrusiu i~jego klockach oraz podsumowanie: \emph{fizyka współczesna nie mówi właściwie, czym jest energia}. Halliday, Resnick i~Walker w~,,Podstawach Fizyki'' (w~wydaniu z~2012 roku) piszą tylko: \emph{Termin energia ma w~rzeczywistości tak szerokie znaczenie, że trudno jest podać jego klarowną definicję}. Następnie przechodzą do operacyjnego wprowadzania rodzajów energii. Hasło \emph{energia} w~,,Słowniku fizycznym'' (Wiedza Powszechna, 1992) rozpoczyna się następująco: \emph{uniwersalna wielkość fizyczna, nadająca się do opisu wszelkiego rodzaju procesów i~oddziaływań występujących w~przyrodzie. Najistotniejszą cechą e. jest to, że podlega ona zasadzie zachowania. E. jest funkcją stanu układu fizycznego.} Wreszcie Wikipedia (dostęp 20.11.2024) podaje: \emph{Energia (gr. $\epsilon\nu\epsilon\rho\gamma\epsilon\iota\alpha$
%ενεργεια
energeia od ἔργον
%{\color{magenta} [tu powinno być to słowo jak w~wikipedii, ale nie wiem jak zrobić tego epsilona z~akcentami]}ἔργον 
ergon „praca”) – skalarna wielkość fizyczna charakteryzująca stan układu fizycznego (materii) jako jego zdolność do wykonania pracy}.

Właśnie tak – jako zdolność do wykonania pracy – najczęściej określa się energię na poziomie podstawowym. Niestety taka definicja cierpi na schorzenie przypominające błędne koło. Podstawowym operacyjnym określeniem energii jest wzór obowiązujący w~mechanice klasycznej i~w~szczególnej teorii względności (podajemy go na poziomie podstawowym):
\begin{equation}\label{en}
\triangle E=\vec{F}\cdot\triangle\vec{r}.
\end{equation}
Przyrost energii ciała (układu ciał) jest równy iloczynowi skalarnemu siły działającej na to ciało i~przemieszczenia (punktu przyłożenia siły). Właśnie iloczyn po prawej stronie nazywamy pracą. Przyrost energii ciała jest równy wykonanej nad tym ciałem pracy. Z~kolei sama energia to zdolność do wykonania pracy. Doprawdy niewielki jest pożytek z~takiego definiowania, najpierw $W$ przez $E$, a~zaraz potem $E$ przez $W$.
\marg{
Podobne błędne koło pojawia się również w~opisie przemian termodynamicznych. Jednak ze względu na niemożliwość obliczenia przyrostów energii w~zderzeniach wszystkich cząstek wprowadzamy makroskopową wielkość zwaną przekazanym ciepłem.
}

Żeby ciało mogło wykonać pracę, musi być zdolne działać siłą. Siła – przypomnijmy -- jest miarą oddziaływania. Pojęcie energii służy do alternatywnego opisu oddziaływania. Każde oddziaływanie można opisać za~pomocą sił lub przekazu energii. Zatem energię lepiej byłoby określić następująco: { \bf energia ciała jest miarą jego zdolności do oddziaływań}. Taka definicja dobrze się sprawdza przy zderzeniach cząstek (nie tylko) elementarnych. Wtedy obserwujemy cząstki powstające w~wyniku oddziaływań. Właśnie energia zderzających się cząstek (w~układzie środka masy) narzuca ograniczenia na to, jakie produkty mogą w~ten sposób powstać. Można ją zatem interpretować jako maksymalną zdolność do oddziaływań.

Określając energię jako miarę zdolności do oddziaływania, usuwamy wspomniane wyżej zapętlenie. Jednak pojawia się problem. Czy zdolność do oddziaływań może być ujemna? Przecież energia potencjalna może być ujemna. Właśnie teraz trzeba przypomnieć znaną, ale rzadko przypominaną prawdę: energia realnych obiektów fizycznych jest zawsze dodatnia.
\marg{
Nie zaglądamy na razie pod horyzont czarnych dziur i~do próżni kwantowej, bo to zbyt nowa fizyka.
}
Wyraźnie wynika to z~najsłynniejszego wzoru nauk ścisłych:
\begin{equation}\label{emc}
  E=mc^2.
\end{equation}
Wzór ten określa energię wewnętrzną\marg{
Patrz tekst \emph{Energia wewnętrzna czy spoczynkowa,} \href{https://www.deltami.edu.pl/2024/06/energia-spoczynkowa-czy-wewnetrzna/}{$\Delta^6\sb{24}$}.
}
-- znacznie częściej zwaną spoczynkową~– każdego ciała. Ta zależność pozwala obliczyć energię całkowitą każdego nieoddziałującego z~otoczeniem ciała w~układzie jego środka masy.  Dotyczy to również układów ciał (każde ciało niebędące cząstką elementarną jest de facto układem ciał). Zatem energia potencjalna oddziaływania składników układu jest tu wliczona.

Zwróćmy też uwagę na to, że zgodnie ze wzorem (\ref{emc}) energia całkowita jest określona jednoznacznie. W~szczególnej teorii względności energia i~pęd tworzą tzw. czterowektor energii-pędu, a~wzór (\ref{emc}) określa jego niezmienniczą długość. Warto pamiętać o~tym bliskim związku energii i~pędu również w~fizyce newtonowskiej, uniknie się wtedy wielu nieporozumień.  Wzór (\ref{en}) warto stosować łącznie z~równaniem określającym zmianę pędu:
\begin{equation}\label{ped}
  \triangle \vec{p}=\vec{F}\cdot\triangle t.
\end{equation}
Z~równań (\ref{en}) i~(\ref{ped}) bezpośrednio wynika także zasada zachowania energii i~pędu. Doprawdy wielkie bogactwo treści kryją w~sobie dwa tak niepozorne i~proste wzory!
Jak zatem pogodzić ze sobą możliwość dowolnego ustalenia wartości energii potencjalnej z~bezwzględną wartością energii wynikającą z~wzoru Einsteina (\ref{emc})? Rozpatrzmy to na przykładzie energii elektrycznej. Każdy ładunek jest źródłem pola elektrycznego, które jest realnym obiektem fizycznym istniejącym również bez towarzystwa ładunków, np. w~fali elektromagnetycznej. \marg{\begin{tikzpicture}
\node at(0,0){\includegraphics[scale=.6]{l1-a}};
\node at(-1.6,-2){\includegraphics[scale=.6]{l1-b}};
\node at(1.6,-2){\includegraphics[scale=.6]{l1-c}};
\draw[thick,->](-.3,-3)--(-.3,-1);
\node[right] at(-.3,-1.3){czas};
\end{tikzpicture}\\[4pt]
  Rys. 1}Pole to ma zawsze energię dodatnią. Przyjrzyjmy się dwóm przeciwnym punktowym ładunkom znajdującym się początkowo nieskończenie daleko od siebie (rys.~1).%
\goodbreak

\marg{
  \begin{tikzpicture}
\node at(0,0){\includegraphics[scale=.6]{l2-a}};
\node at(-1.6,-2){\includegraphics[scale=.6]{l2-b}};
\node at(1.6,-2){\includegraphics[scale=.6]{l2-c}};
\draw[thick,->](-.3,-3)--(-.3,-1.1);
\node[right] at(-.3,-1.4){czas};
\end{tikzpicture}\\[4pt]
  Rys. 2}Jeśli poczekamy dostatecznie długo, zobaczymy, że ładunki zbliżają się do siebie wskutek kulombowskiego przyciągania. Ponieważ rośnie ich energia kinetyczna, maleje energia wspólnego pola elektrycznego. Gdy ładunki się zetkną, pole całkowicie zniknie. Ponieważ pole nie jest widoczne, a~poruszające się ładunki zaobserwować bardzo łatwo, fizycy właśnie im przypisali niewidoczną energię pola. Nazwali ją energią potencjalną dwóch ładunków i~dla wygody ustalili, że dla ładunków nieskończenie od siebie odległych jest ona równa zeru, tak jak siła między nimi. Niestety w~rozważanym przypadku energia pola jest wtedy największa! Zatem tak określona energia potencjalna dwóch różnoimiennych ładunków jest zawsze ujemna (w~granicy równa zeru). Jednak energia ich pola jest zawsze dodatnia (równa zeru tylko wtedy, gdy pole całkowicie znika).

Przyjrzyjmy się temu trochę bliżej. Możemy stwierdzić, że praca wykonana przez siłę elektryczną (równa energii kinetycznej ładunków tuż przed ich anihilacją) jest równa początkowej energii pola elektrycznego. Jeśli energię pola pochodzącego od pojedynczego ładunku $q$ oznaczymy jako $E(q)$, to energią pola dwóch znajdujących się nieskończenie daleko od siebie ładunków będzie $2E(q)$, możemy zatem napisać:
\marg{
Energia pola ładunku punktowego okazuje się nieskończenie wielka. Ten problem rozwiązuje się w~kwantowej teorii pola dzięki tzw. procedurze renormalizacji.
}
\begin{equation}\label{praca}
  W=2E(q).
\end{equation}
Rozważmy teraz inny proces. Niech dwa ładunki dodatnie o~wartości $q$ będą nieskończenie daleko od siebie (rys. 2).


Jeśli chcemy je połączyć w~jeden ładunek o~wartości $2q$, to teraz my musimy wykonać pracę o~tej samej wartości co poprzednio. Wynikiem tego procesu jest jedno pole ładunku $2q$ powstałe z~początkowych dwóch pól dzięki wykonanej pracy. Zatem:
\begin{equation}\label{eew}
  E(2q)=2E(q)+W.
\end{equation}
Wstawiając do równania (\ref{eew}) wzór (\ref{praca}), otrzymujemy:
$$
E(2q)=4E(q).
$$
Pole ładunku $2q$ różni się od pola ładunku $q$ tylko natężeniem w~każdym punkcie, a~powyższy wzór mówi, że energia pola pochodzącego od ładunku~$2q$ jest czterokrotnie większa od energii pola ładunku $q$. Stąd możemy wywnioskować, że gęstość energii pola elektrycznego jest wprost proporcjonalna do kwadratu natężenia. Konkretny wzór określający gęstość tej energii łatwo otrzymujemy dla kondensatora płaskiego, bo tam pole jest jednorodne i~zajmuje ograniczony obszar.

\vspace*{\parskip}
\begin{dwieszpalty}Warto przypomnieć inną ważną cechę energii – jest ona funkcją stanu. Niestety wielu o~tym nie pamięta. Paul G. Hewitt w~znakomitym kursie podstawowym fizyki (,,Fizyka wokół nas'') przedstawia energię w~następujący sposób: \emph{Chociaż z~pojęciem tym spotykamy się na co dzień, to jednak energię trudno zdefiniować. Nie jest ona rzeczą, lecz łączy w~sobie pojęcie rzeczy z~pojęciem procesu, tak jakby była jednocześnie i~rzeczownikiem, i~czasownikiem. (\dots) Sama materia jest swoistym koncentratem energii, o~czym świadczy słynny wzór Einsteina\ldots} Takie ,,mistyczne'' traktowanie energii i~mieszanie tego pojęcia z~pojęciem materii jest – niestety – mocno rozpowszechnione. Co chwila można się o~taki mistycyzm potknąć przy lekturze tekstów popularnonaukowych, nieraz pisanych przez wybitnych fizyków. Nie warto cytować licznych przykładów, podajmy tylko jeden – powszechnie stosowany zapis reakcji fotosyntezy:
$$
6\mathrm{H_2O+6CO_2}+h\nu\, \mbox{(energia świetlna)}\rightarrow \mathrm{C_6H_{12}O_6+6 O_2}\!\!\uparrow.
$$
Ten pochodzi z~Wikipedii. Co tu się nie zgadza? Otóż w~reakcji uczestniczą różne cząsteczki oraz\ldots\ energia świetlna!  Obok realnych cząsteczek w~reakcji uczestniczy funkcja stanu nazywana energią. A~przecież wystarczy zamiast ,,energia świetlna'' napisać ,,fotony'', będące tak samo realnymi obiektami jak elektrony czy atomy.

Zmierzamy już do końca, więc podsumujmy. Energia to funkcja stanu układu fizycznego określająca jego zdolność do oddziaływań z~innymi układami. W~skrócie – energia jest miarą zdolności do oddziaływań. Trochę to podobne do popularnej „intuicyjnej” definicji masy – to miara bezwładności ciała. Czemu jednak energia jest tak ściśle związana z~masą słynnym wzorem Einsteina? Okazało się, że masa wszystkich cząstek elementarnych pochodzi z~ich oddziaływania z~polem Higgsa. Więc może nie jest tak tajemnicze to, że zdolność cząstki do oddziaływania (z~polem Higgsa) jest równocześnie miarą jej bezwładności?

Mam nadzieję, że Kip S. Thorne by się w~tym momencie uśmiechnął. Bo dobra definicja wielkości fizycznej pomaga zrozumieć jej istotne cechy i~może zapoczątkować ,,owocne badania''.




\end{dwieszpalty}
